\documentclass[11pt]{ctexart}
\usepackage[a4paper,margin=1in]{geometry}
\usepackage{graphicx}
\usepackage{xcolor}
\usepackage{hyperref}
\definecolor{refgray}{gray}{0.45}
\title{\textbf{Market Views｜全球投行叙事汇编}\\[0.4em]{\large AI硬件封装瓶颈与开源模型扩张并行，中国出口结汇与去库分化驱动外汇与商品结构性行情}}
\author{KC桌面 自动生成}
\date{260723}
\begin{document}
\maketitle
\tableofcontents
\newpage
\section{一页摘要}
\begin{itemize}
  \item 亚洲市场去杠杆进程过半但波动率未正常化，中国出口企业结汇意愿回升支撑人民币，政策层对升值容忍度提高形成正反馈。
  \item AI基础设施投资重塑台湾、越南、马来西亚增长与通胀动态，部分央行可能转向更紧货币政策。
  \item 全球贸易增长高度依赖AI商品，非AI贸易停滞，美加贸易摩擦与韩国出口走弱加剧外汇不确定性。
  \item AI软件变现仍处早期，安全与SaaS收入拐点需2-3年，企业部署多处于沙盒阶段。
  \item 开源模型参数规模持续增长，每个下载实例创造独立HBM/DRAM需求，存储TAM不降反升；CoWoS-L面临9.5倍光罩尺寸焊接可靠性瓶颈，可能制约下一代AI芯片封装密度。
  \item 中国工业与材料板块分化显著：电池材料受益储能高景气与锂供给约束，铜铝库存去化加速，AI数据中心电源订单积压延伸至2027年。
  \item 中国消费与房地产分化加剧，龙头韧性凸显：房地产销售边际改善但城市间差异大，高端白酒定价能力凸显，黄金珠宝需求结构性改善。
  \item 全球权益市场正从少数头部公司向更广泛板块扩散，动量因子出现极端反转。
\end{itemize}
\section{全球投行叙事汇编}
今日纳入 60 篇投行/券商研究，按 9 个市场、资产与产业主题系统整理。
\newpage
\subsection{亚洲市场：去杠杆、结汇与AI驱动的结构性变化}
\textbf{亚洲市场正经历多重结构性调整：韩国去杠杆进程过半但波动率未正常化，中国出口企业结汇意愿回升支撑人民币，AI基础设施投资重塑部分经济体的增长与通胀动态，而全球LNG市场面临供应中断与低库存的短期压力。}
\subsubsection*{主流共识}
\begin{itemize}
  \item 韩国市场去杠杆已部分完成，但波动率指标尚未回归正常，市场自我修正而非趋势逆转。
  \item 中国出口企业结汇意愿回升，政策层对人民币升值的容忍度提高，形成正反馈循环。
  \item AI基础设施投资正在显著提升台湾、越南、马来西亚等地区的经济增长，并可能推动部分央行转向更紧的货币政策。
\end{itemize}
\subsubsection*{分歧与条件差异}
\begin{itemize}
  \item 韩国去杠杆的后续影响：JPM认为基本面仍有支撑，但波动率持续高位对资产管理构成约束。
  \item 人民币升值路径：DB强调结汇驱动，但资本外流恢复可能部分对冲升值压力。
  \item LNG价格走势：MS认为供应中断和低库存支撑价格，但地缘缓和可能带来调整。
\end{itemize}
\subsubsection*{机构观点}
\paragraph{MS} 全球LNG市场面临低库存、供应中断和季节性需求三重叠加，价格上行不确定性未出清。亚洲JKM现货价格7月以来上涨35\%，年内涨幅超一倍。卡塔尔出口恢复延迟，Golden Pass项目调试波动，Freeport LNG检修至8月下旬。美国LNG出口商Venture Global和Cheniere Energy报价仍低于资产净值，分别有约21\%和4\%的上行空间。2026年全球LNG预计出现约1300万吨缺口。
{\small\color{refgray}数据：JKM价格7月以来上涨35\%，年内涨幅超一倍\par}
{\small\color{refgray}数据：全球液化装置利用率7月达95\%，高于五年均值83\%\par}
{\small\color{refgray}数据：Venture Global估值较当前报价高出约21\%\par}
{\small\color{refgray}数据：2026年全球LNG预计出现约1300万吨缺口\par}
{\small\color{refgray}边际变化：从关注中东地缘事件转向卡塔尔出口恢复速度和Golden Pass爬坡对短期供需平衡的决定性影响。\par}
\paragraph{DB} 6月银行代客净美元卖盘升至574亿美元，为2月以来新高，出口企业结汇比率从58.2\%反弹至64.5\%。美元/人民币中间价持续走低，表明政策层对人民币升值的容忍度上升，提高出口企业持有美元的机会成本。即使结汇比率稳定，贸易顺差扩张也将增加可结汇规模。资本外流恢复（证券投资从净流入180亿转为净流出360亿）可能部分对冲升值压力，但整体跨境资金流仍支撑人民币。
{\small\color{refgray}数据：6月银行代客净美元卖盘574亿美元，为2月以来最高\par}
{\small\color{refgray}数据：出口企业结汇比率从58.2\%反弹至64.5\%\par}
{\small\color{refgray}数据：证券投资从净流入180亿美元转为净流出360亿美元\par}
{\small\color{refgray}数据：境外机构增持中国国债约10亿美元\par}
{\small\color{refgray}边际变化：从关注结汇比率反弹转向政策中间价信号与结汇行为形成的正反馈循环，以及资本外流恢复对升值压力的对冲。\par}
\paragraph{HSBC} AI对亚洲的影响不再是远期叙事，正在重塑通胀动态和央行政策权衡。AI硬件出口在台湾和越南已达到相当于GDP近一半的水平，数据中心投资在马来西亚尤为突出。美国四大云服务商AI资本支出预计从2022年约1500亿美元增至2026年超过7000亿美元。上游成本不确定性通过技术供应链传导至更广泛价格体系，可能推动韩国、台湾等央行转向更紧货币政策。
{\small\color{refgray}数据：AI硬件出口在台湾和越南相当于GDP近一半\par}
{\small\color{refgray}数据：美国四大云服务商AI资本支出预计2026年超7000亿美元\par}
{\small\color{refgray}数据：台积电全球晶圆代工市占率约70\%\par}
{\small\color{refgray}数据：三星和SK海力士合计占全球DRAM市场约70\%\par}
{\small\color{refgray}边际变化：从关注AI对增长的远期叙事转向AI基础设施投资对通胀动态和央行政策权衡的即时影响。\par}
\paragraph{JPM} 韩国KOSPI指数自6月22日高点下跌约28\%，杠杆ETF和hedge fund仓位平仓放大跌幅。去杠杆已部分完成：杠杆ETF规模从500亿美元降至260亿美元（完成约75\%），多空比率从5.5倍降至4倍以下（完成超50\%）。但VKOSPI与VIX比值仍接近5倍（历史常态约1倍），波动率未正常化。外资流出超1100亿美元，90\%来自两只存储股。基本面仍有支撑，去杠杆是市场自我修正。
{\small\color{refgray}数据：KOSPI自6月22日高点下跌约28\%\par}
{\small\color{refgray}数据：杠杆ETF规模从500亿美元降至260亿美元，去杠杆完成约75\%\par}
{\small\color{refgray}数据：多空比率从5.5倍降至4倍以下，去杠杆完成超50\%\par}
{\small\color{refgray}数据：VKOSPI与VIX比值接近5倍，历史常态约1倍\par}
{\small\color{refgray}边际变化：从关注韩国市场下跌本身转向去杠杆进程的结构性评估，强调波动率指标尚未正常化对资产管理的持续约束。\par}
\paragraph{DB} 百胜中国将于7月30日发布二季度业绩。DB在最新报告中预计，公司二季度总收入与经营利润将分别实现中个位数和高个位数同比增长，净利润约在2.3至2.4亿美元区间。这份报告的核心看点并非业绩数字本身，而是两个将在业绩会上被重点讨论的议题：12亿美元必胜客中国品牌收购的债务融资安排，以及6月下旬以来南方洪涝与北方强降水对旺季销售的实际影响。
{\small\color{refgray}数据：报告维持对百胜中国的正面看法。当前报价约347.8港元，对应2026年预期市盈率约15.5倍，低于其三年历史均值超过一个标准差。该机构注意到，在AI研究热潮中，境内外读者对中国消费股表现出对低估值、高股东回报标的的偏好，百胜中国约10\%的股东回报率使其在近期有所反弹。\par}
{\small\color{refgray}数据：DB预计二季度餐厅利润率同比微降0.3个百分点，主要不确定性来自人力成本。外卖占比在2026年上半年仍可能继续上升，这直接推高了劳动力支出。但报告同时指出，公司通过持续谈判降低新店和现有门店的租金成本，在租金及其他成本项上持续节省，这部分对冲了人力端的不确定性。\par}
\subsubsection*{关键数据}
\begin{itemize}
  \item JKM价格7月以来上涨35\%，年内涨幅超一倍
  \item 6月银行代客净美元卖盘574亿美元，为2月以来最高
  \item 出口企业结汇比率从58.2\%反弹至64.5\%
  \item AI硬件出口在台湾和越南相当于GDP近一半
  \item 美国四大云服务商AI资本支出预计2026年超7000亿美元
  \item 韩国杠杆ETF规模从500亿美元降至260亿美元，去杠杆完成约75\%
  \item VKOSPI与VIX比值接近5倍，历史常态约1倍
  \item 2026年全球LNG预计出现约1300万吨缺口
\end{itemize}
\begin{figure}[htbp]
\centering
\includegraphics[width=0.88\linewidth]{figures/fig_027.jpg}
\caption{Exhibit 1 - Exhibit 1: After a dip in June following MoU news, share prices have begun to recover for LNG stocks in our coverage. Exhibit 2: JKM prices have rallied, in-line with our constructive view. Our 2026 forecasts remain modestly abov}
\end{figure}
\begin{figure}[htbp]
\centering
\includegraphics[width=0.88\linewidth]{figures/fig_050.jpg}
\caption{Figure 1 - Figure 1: Net USD selling rebounds as exporters' FX conversion gains traction Figure 2: Exporters' FX conversion has rebounded despite a still-weak 3MMA \#\# Appendix 1}
\end{figure}
\begin{figure}[htbp]
\centering
\includegraphics[width=0.88\linewidth]{figures/fig_102.jpg}
\caption{Figure 1 - Figure 1: KOSPI index has declined \$-29\textbackslash{}\%\$ from the peak of 22 Jun Figure 2: Bringing the index RSI to lowest levels since the start of this rally post Liberation Day in 2025 Figure 3: Korea is still the best performing market}
\end{figure}
\begin{figure}[htbp]
\centering
\includegraphics[width=0.88\linewidth]{figures/fig_103.jpg}
\caption{Figure 1 - Figure 4: KOSPI has largely been a reflection of memory upside in recent months}
\end{figure}
{\small\color{refgray}本节综合 5 篇研究；涉及 MS、DB、HSBC、JPM。}
\newpage
\subsection{全球贸易与外汇策略：AI驱动分化与政策不确定性}
\textbf{全球贸易增长高度依赖AI相关商品，非AI贸易增长停滞，同时美加贸易摩擦与韩国出口数据走弱加剧了外汇市场的不确定性；量化策略与主观判断的分歧凸显了美联储政策路径的关键作用。}
\subsubsection*{主流共识}
\begin{itemize}
  \item AI相关商品是当前全球贸易增长的核心驱动力，贡献了2025年全球出口增长的80\%以上。
  \item 非AI商品贸易增长自2024年以来基本持平，贸易增长呈现结构性分化。
  \item 美加贸易摩擦升级将抑制加拿大经济增长，增加加拿大央行政策不确定性。
\end{itemize}
\subsubsection*{分歧与条件差异}
\begin{itemize}
  \item MS量化策略显示温和美元多头，但策略团队基于美联储更宽松预期持中性偏空立场。
  \item 韩国7月出口数据走弱，半导体出口持平但其他品类下滑，市场对全球贸易复苏广度存在分歧。
  \item 美加关税是否全面实施及加拿大是否报复，将影响通胀传导速度和央行政策路径。
\end{itemize}
\subsubsection*{机构观点}
\paragraph{MS} MS量化多因子策略当前持有温和美元多头仓位（净敞口-21.0\%），但策略团队自身对美元持中性偏空立场，核心依据是该行对美联储政策路径的预期比市场定价更为宽松。报告认为，量化策略未能充分捕捉美联储政策差异这一宏观变量，因此两者分歧并非系统性问题，而是输入条件不同。该策略夏普比率1.87，年化回报9.3\%，最大回撤-6.8\%，偏好利差交易。
{\small\color{refgray}数据：策略净美元敞口-21.0\%\par}
{\small\color{refgray}数据：最大空头头寸离岸人民币权重-47.2\%\par}
{\small\color{refgray}数据：夏普比率1.87，年化回报9.3\%，最大回撤-6.8\%\par}
{\small\color{refgray}边际变化：量化信号与主观判断的分歧持续存在，但报告强调量化策略可作为压力测试工具，而非直接驱动观点。\par}
\paragraph{GS} GS估算美国对加拿大200亿美元商品加征关税将在未来12个月内对加拿大实际GDP增长造成0.2-0.3个百分点的影响，主要来自出口增长节奏变化（0.2pp）和不确定性对投资的抑制（0.2pp），部分被消费者转向本土替代品所抵消。在加拿大有限报复的基准假设下，通胀上行风险有限，GS维持加拿大央行2026年政策利率2.25\%不变的预测。
{\small\color{refgray}数据：关税覆盖200亿美元加拿大商品\par}
{\small\color{refgray}数据：有效关税率提高约2.5个百分点\par}
{\small\color{refgray}数据：GDP影响0.2-0.3个百分点\par}
{\small\color{refgray}边际变化：关税宣布后30天窗口期增加政策不确定性，但加拿大未采取报复措施，与去年态度形成对比。\par}
\paragraph{GS} 韩国7月前20天日均出口额环比下降2.3\%，结束此前两个月增长。半导体出口环比仅降0.3\%基本持平，但其他科技产品出口环比下降11.2\%（计算机降23.1\%），非科技产品出口下降5.6\%。对美国和台湾出口降幅占整体下滑的三分之二，对中国大陆出口则连续9个月扩张。贸易顺差从6月历史高点174亿美元收窄至122亿美元。
{\small\color{refgray}数据：日均出口环比-2.3\%\par}
{\small\color{refgray}数据：半导体出口环比-0.3\%\par}
{\small\color{refgray}数据：其他科技产品出口环比-11.2\%\par}
{\small\color{refgray}数据：对美出口环比-10.4\%，对台出口环比-28.4\%\par}
{\small\color{refgray}边际变化：此前两个月半导体与其他品类同步扩张，7月出现分化，半导体持平而其他品类下滑，显示贸易复苏广度收窄。\par}
\paragraph{HSBC} HSBC报告指出，AI相关商品在2025年贡献全球商品贸易增长超过40\%，2026年第一季度进一步攀升至约80\%，尽管仅占贸易商品总量的约六分之一。非AI商品出口自2024年以来基本持平。亚洲是AI出口核心，中国大陆、台湾和香港合计占全球AI商品出口总值的44\%，台湾约80\%总出口与AI价值链相关。若AI投资周期降温，全球出口增长每年可能减少0.2-0.3个百分点。
{\small\color{refgray}数据：AI商品贡献2025年全球贸易增长>40\%\par}
{\small\color{refgray}数据：2026年Q1 AI商品贡献全球出口增长约80\%\par}
{\small\color{refgray}数据：AI商品占全球贸易商品总量约1/6\par}
{\small\color{refgray}数据：中国大陆、台湾、香港合计占全球AI出口44\%\par}
{\small\color{refgray}边际变化：AI商品在全球商品贸易中的占比从2024年平均14\%上升至当前近20\%，非AI贸易增长停滞，贸易增长对AI周期的依赖度显著提升。\par}
\paragraph{GS} 中国家电行业正经历一个关键的分化窗口。相关研究在最新发布的6月家电追踪报告中指出，国内零售与出货数据在低基数效应下出现环比改善，而出口端则凭借中国企业的份额扩张保持了超出预期的韧性。这份报告的核心判断是：国内需求最剧烈的调整阶段可能已经过去，但成本传导能力在需求弱复苏背景下仍然受限；出口的持续增长更多来自产品创新与全球竞争力，而非单纯的周期性需求回暖。
{\small\color{refgray}数据：6月空调国内出货量同比下降7\%，较5月的-15\%明显收窄，也显著好于此前产业在线（IOL）生产计划所预示的-23\%降幅。研究认为，这一改善部分得益于618促销期间空调折扣力度加大，拉动了终端需求。统计局数据显示，6月家电和电子产品零售额同比下降9\%，较5月的-16\%有所回升。报告预计，随着基数效应进一步走弱，三季度国内零售增速有望转正至同比增长15\%，四季度...\par}
{\small\color{refgray}数据：观察评论： 国内出货数据改善更多是基数效应和促销驱动的短期反弹，而非需求端的结构性复苏。报告特别指出，在缺乏强劲需求恢复的情况下，企业难以通过提价完全转嫁成本不确定性——这意味着二季度以来原材料成本上涨对毛利率的侵蚀，可能在三季度仍会持续体现。\par}
\subsubsection*{关键数据}
\begin{itemize}
  \item MS量化策略净美元敞口-21.0\%，最大空头CNH权重-47.2\%
  \item 美对加200亿美元商品关税，GDP影响0.2-0.3个百分点
  \item 韩国7月前20天日均出口环比-2.3\%，半导体环比-0.3\%，其他科技-11.2\%
  \item 韩国贸易顺差从174亿收窄至122亿美元
  \item AI商品贡献2026年Q1全球出口增长约80\%，占贸易总量约1/6
  \item 台湾约80\%总出口与AI价值链相关
  \item 中国大陆对韩国出口连续9个月扩张
\end{itemize}
\begin{figure}[htbp]
\centering
\includegraphics[width=0.88\linewidth]{figures/fig_039.jpg}
\caption{Exhibit 1 - Exhibit 1: Strategy has a modest long USD position Exhibit 2: Which are the favourite longs and shorts? \textbackslash{}* Why are we not long USD, if the systematic strategy is? The strategy's net long USD position is sending a different sign}
\end{figure}
\begin{figure}[htbp]
\centering
\includegraphics[width=0.88\linewidth]{figures/fig_068.jpg}
\caption{Exhibit 1 - Exhibit 1: The White House Announced Tariffs on \textbackslash{}\$20bn of Canadian Imports Joseph Briggs \#\# Disclosure Appendix \#\# Reg AC}
\end{figure}
\begin{figure}[htbp]
\centering
\includegraphics[width=0.88\linewidth]{figures/fig_086.jpg}
\caption{Exhibit 1 - Exhibit 1: Pullback in early July exports \#\# Disclosure Appendix \#\# Reg AC We, Irene Choi and Goohoon Kwon, CFA, hereby certify that all of the views expressed in this report accurately reflect our personal views, which have not}
\end{figure}
\begin{figure}[htbp]
\centering
\includegraphics[width=0.88\linewidth]{figures/fig_038.jpg}
\caption{Exhibit 2 - Exhibit 1: Strategy has a modest long USD position Exhibit 2: Which are the favourite longs and shorts? \textbackslash{}* Why are we not long USD, if the systematic strategy is? The strategy's net long USD position is sending a different sign}
\end{figure}
{\small\color{refgray}本节综合 5 篇研究；涉及 MS、GS、HSBC。}
\newpage
\subsection{AI变现：从叙事到数字的过渡期}
\textbf{AI软件变现仍处早期，安全与SaaS收入拐点需2-3年；中国模型智能差距缩小但成本优势显著，竞争焦点转向平台与系统级能力。}
\subsubsection*{主流共识}
\begin{itemize}
  \item AI变现仍处早期，SaaS公司AI收入占比过小，不足以触发盈利预期上调。
  \item 企业AI部署多处于沙盒阶段，安全预算拐点预计在2027-2028年。
  \item 中国AI模型迭代加速，开源策略扩大生态但牺牲短期变现。
  \item 数据云需求上升，Databricks与Snowflake均受益，但增速分化明显。
\end{itemize}
\subsubsection*{分歧与条件差异}
\begin{itemize}
  \item AI安全支出拐点时间：GS认为最早2026Q4-2027H1，而市场已提前定价中期拐点。
  \item 中国模型竞争力：部分基准上中国模型首次领先美国，但成本效率优势是否可持续存在分歧。
  \item EDA软件护城河：Citi认为AI难以颠覆，但Kimi K3绕开EDA的案例引发短期冲击。
  \item AI眼镜规模化：JEF看好Meta先发优势，但隐私监管与续航问题可能限制渗透率。
\end{itemize}
\subsubsection*{机构观点}
\paragraph{BARC} 美国软件Q2财报季不太可能成为AI变现转折点。SaaS公司AI收入占比太小，无法推动盈利预期上调；基础设施类公司因AI部署需求延续增长。空头仓位上升反映市场对AI变现的怀疑。
{\small\color{refgray}数据：HubSpot空头比例从4月7.1\%升至7月12.4\%\par}
{\small\color{refgray}数据：ServiceNow空头比例从3.8\%升至6.0\%\par}
{\small\color{refgray}数据：monday.com空头比例从14.8\%升至20.2\%\par}
{\small\color{refgray}边际变化：空头仓位上升，市场对AI变现怀疑正在定价。\par}
\paragraph{Citi} 中国AI市场从模型竞赛转向平台竞争。Qwen3.8 Max与Kimi K3密集发布催生“模型无关”采购环境，企业采用多模型策略，竞争焦点转向托管与管理多个模型的平台。EDA软件护城河深，AI代理更可能成为工具而非替代者。
{\small\color{refgray}数据：Qwen3.8 Max参数2.4万亿，Kimi K3参数2.8万亿\par}
{\small\color{refgray}数据：西门子EDA年收入约21亿欧元，全球份额15\%\par}
{\small\color{refgray}数据：西门子估值折价近期扩大，SOTP显示EDA价值未被充分反映\par}
{\small\color{refgray}边际变化：中国AI迭代速度超出预期，竞争焦点从模型转向平台。\par}
\paragraph{GS} AI安全支出拐点最早2026Q4-2027H1，参考云安全周期约5年滞后。当前安全预算未实质变化，报价已提前反映中期拐点。AI安全收入2025年仅1-2亿美元，渗透率0.6\%。中国软件WAIC新品密集发布，但金山办公AI付费转化慢于预期，评级下调。
{\small\color{refgray}数据：2025年AI安全收入1-2亿美元，对应AI推理IaaS收入170亿美元\par}
{\small\color{refgray}数据：AI安全预算拐点最早2026Q4-2027H1\par}
{\small\color{refgray}数据：金山办公目标价从326元降至201元，隐含PE从45倍降至31倍\par}
{\small\color{refgray}边际变化：安全预算尚未变化，报价已提前定价中期拐点；金山办公AI付费转化慢于预期。\par}
\paragraph{JPM} 中国AI芯片竞争基线重置，Hopper级单芯片性能成入场门槛。Prefill/Decode分离架构成行业标准，超节点方案应对制程瓶颈与万亿参数模型。AI算力从能力建设转向Token经济，Agentic AI创造指数级需求。
{\small\color{refgray}数据：沐曦Tiangai 300推理性能比Hopper高10-20\%\par}
{\small\color{refgray}数据：Kimi K3参数2.8万亿\par}
{\small\color{refgray}数据：华为Atlas 950 SuperPoD功耗1.7MW，是英伟达Rubin NVL72的7.4倍\par}
{\small\color{refgray}边际变化：单芯片性能竞赛转向系统级部署，超节点方案成关键架构。\par}
\paragraph{JEF} 中国AI策略转向务实，开源模式向新兴市场输出算力与绿色能源。K3模型在编码基准上首次领先美国，成本效率优势显著。Databricks融资估值1880亿美元，为Snowflake提供参照。Meta AI眼镜有望成为首个规模化出货的轻量AI交互界面。
{\small\color{refgray}数据：K3在Frontend Code Arena首次领先美国模型\par}
{\small\color{refgray}数据：中国模型API混合价格仅为美国模型的几分之一\par}
{\small\color{refgray}数据：Databricks估值1880亿美元，较1340亿提升40\%\par}
{\small\color{refgray}数据：Meta AI眼镜硬件收入机会140-180亿美元\par}
{\small\color{refgray}边际变化：中国模型智能差距缩小，成本效率优势显著；AI眼镜规模化出货预期上升。\par}
\subsubsection*{关键数据}
\begin{itemize}
  \item HubSpot空头比例从7.1\%升至12.4\%
  \item 2025年AI安全收入1-2亿美元，渗透率0.6\%
  \item K3在Frontend Code Arena首次领先美国模型
  \item 中国模型API混合价格仅为美国模型的几分之一
  \item Databricks估值1880亿美元，较1340亿提升40\%
  \item Meta AI眼镜硬件收入机会140-180亿美元
  \item 华为Atlas 950 SuperPoD功耗1.7MW，是英伟达Rubin NVL72的7.4倍
  \item 金山办公目标价从326元降至201元
\end{itemize}
\begin{figure}[htbp]
\centering
\includegraphics[width=0.88\linewidth]{figures/fig_001.jpg}
\caption{FIGURE 1 - \textbackslash{}- ServiceNow: We maintain our Overweight rating and price target of \textbackslash{}\$134 based on 20x EV/ CY27E FCF of \textbackslash{}\$6.91bn (both unchanged). We update our model primarily to account for updated FX impacts. \textbackslash{}- WhiteFiber: We maintain our Equal Weight rating and raise ou}
\end{figure}
\begin{figure}[htbp]
\centering
\includegraphics[width=0.88\linewidth]{figures/fig_049.jpg}
\caption{Figure 1 - \#\# CITI'S TAKE EDA software (Electronic Design Automation) in context - We estimate that Siemens has about 15\% market share in EDA, behind Synopsys and Cadence. Siemens's \textbackslash{}\textasciitilde{}€2.1bn EDA revenue is equivalent to \textbackslash{}\textasciitilde{}11\% of DI divisional sales, or 4\% of group sales }
\end{figure}
\begin{figure}[htbp]
\centering
\includegraphics[width=0.88\linewidth]{figures/fig_052.jpg}
\caption{Exhibit 1 - Exhibit 3: Security fundamentals have yet to inflect higher; look for improving lead indicators (e.g. bookings) into 2HCY}
\end{figure}
\begin{figure}[htbp]
\centering
\includegraphics[width=0.88\linewidth]{figures/fig_120.jpg}
\caption{Exhibit 1 - Exhibit 1 - META AI Glasses Are a Multi-Billion-Dollar Opportunity ASP= Average selling price}
\end{figure}
{\small\color{refgray}本节综合 10 篇研究；涉及 BARC、Citi、GS、JPM、JEF。}
\newpage
\subsection{AI硬件与半导体：开源模型扩张与封装瓶颈下的结构性需求}
\textbf{开源大模型权重规模持续增长，每个下载实例都产生独立的HBM/DRAM需求，存储芯片TAM不降反升；同时先进封装（CoWoS-L）面临9.5倍光罩尺寸的焊接可靠性瓶颈，可能制约下一代AI芯片的封装密度提升节奏。企业IT支出在AI基础设施、安全需求和本地部署回流驱动下呈现结构性增长，订单积压已延伸至2027年。}
\subsubsection*{主流共识}
\begin{itemize}
  \item 开源模型参数规模持续增长（如Kimi K3达2.8T参数），MoE架构下全部专家权重必须常驻HBM，存储容量与总参数成正比，每个下载实例都创造独立的内存需求。
  \item 企业IT支出正经历由AI基础设施、安全需求和本地部署回流共同驱动的结构性增长，收入增长同时来自出货量和价格，订单积压已延伸至2026年下半年和2027年。
  \item 先进封装中基板刚性、平坦度与热膨胀系数匹配成为决定下一代封装量产的关键约束，日本材料与基板厂商拥有难以替代的竞争优势。
  \item AI PCB供应商长期份额取决于高质量制造产能而非初期份额分配，产能扩张速度和良率爬坡能力比第一轮分配更能决定长期地位。
\end{itemize}
\subsubsection*{分歧与条件差异}
\begin{itemize}
  \item CoWoS-L能否在2027年前解决9.5倍光罩尺寸的焊接问题存在分歧：JPM研讨会指出5.5倍尺寸的Rubin Ultra已遇到困难，9.5倍问题仍待解决；而英特尔EMIB-T路径虽可支持12倍光罩尺寸，但量产能力尚未建立。
  \item 康宁二季度营收超预期幅度分歧：MS测算超预期空间在5000万美元左右，不太可能超过1亿美元，但市场共识已隐含约2亿美元环比增长，实际结果若显著高于预期可能来自定价因素而非销量。
  \item AI PCB供应商VGT在Vera Rubin项目中的初期份额低于上一代，但JPM认为这不是结构性份额流失，随着产能扩大将重新获得份额；市场短期博弈初期分配，但长期格局取决于产能支撑能力。
\end{itemize}
\subsubsection*{机构观点}
\paragraph{BofA} 开源模型权重规模持续增长，从DeepSeek-V3的671B参数（2024年12月）到Kimi K3的2.8T参数（2026年7月），两年内增长超4倍。即使采用MXFP4低比特量化将存储需求压缩一半，K3所需权重内存仍达1.4TB，是DeepSeek-V3在FP8精度下的两倍。MoE架构下全部专家权重必须常驻HBM，存储容量与总参数成正比。每个开源模型下载实例都创造独立的客户侧内存需求，这与闭源API模式下共享内存不同，开源模型的扩散反而扩大了存储芯片TAM。
{\small\color{refgray}数据：Kimi K3参数规模2.8T，是DeepSeek-V3（671B）的4倍以上\par}
{\small\color{refgray}数据：K3在MXFP4下仍需1.4TB HBM，是DeepSeek-V3 FP8下的两倍\par}
{\small\color{refgray}数据：开源模型API价格仅为Claude Opus的1/5，但每个服务实例仍需1.4TB HBM\par}
{\small\color{refgray}边际变化：市场此前认为开源模型效率提升会降低存储需求，但BofA指出参数规模增长速度远超量化压缩节省，且每个下载实例都创造独立内存需求，存储TAM反而扩大。\par}
\paragraph{MS} 康宁二季度营收超预期幅度有限（约5000万美元，不超过1亿美元），市场关注点转向利润率表现和光学业务在三季度的供应驱动型加速。光学业务处于售罄状态，新增产能上线后贡献至少5000万美元超预期增量。太阳能业务仍是最大不确定性，二季度对每股收益影响约0.07美元（含0.03美元临时停产费用和0.04美元硅片爬坡低效）。自6月底GlassBridge消息后报价回调约40\%，但估值倍数（约35倍2028年每股收益）远高于过去十年16倍均值，有意义的报价上行需看到额外定价证据或更快的供应爬坡。
{\small\color{refgray}数据：二季度营收超预期幅度约5000万美元，不超过1亿美元\par}
{\small\color{refgray}数据：太阳能业务对每股收益影响约0.07美元，其中0.03美元为停产费用\par}
{\small\color{refgray}数据：报价自高点回调约40\%，目标价180美元对应35倍2028年每股收益\par}
{\small\color{refgray}边际变化：市场此前关注二季度营收超预期，但MS指出售罄状态限制了弹性，焦点转向三季度供应驱动型加速和利润率表现。\par}
\paragraph{JPM} 先进封装面临CoWoS-L 9.5倍光罩尺寸的焊接可靠性瓶颈，5.5倍尺寸的Rubin Ultra已遇到困难，9.5倍问题在2027年上半年仍待解决。若无法成功，英伟达Feynman世代（2029年）的实现将直接受限。英特尔EMIB-T路径可支持12倍光罩尺寸且翘曲问题较小，但IFS对外量产能力尚未建立。日本Ibiden等厂商在大尺寸、高刚性、高平坦度基板领域竞争优势凸显。企业IT支出方面，订单积压已延伸至2027年，内存成本上升导致OEM报价有效期从60-90天缩短至1-2周，存储和AI集群交付周期超6个月。Dell在AI工厂建设中表现突出，存储端AI导向厂商（Vast、Weka）势头较猛。AI PCB供应商VGT在Vera Rubin初期份额低于上一代，但产能扩张将使其在后续大批量生产中重新获得份额；Kyber机架延迟的影响可被当前代际强劲需求和谷歌项目份额扩大部分抵消。
{\small\color{refgray}数据：CoWoS-L光罩尺寸从3.3倍扩大到9.5倍，基板尺寸从85mm见方增至130mm x 140mm\par}
{\small\color{refgray}数据：OEM报价有效期从60-90天缩短至1-2周，部分按发货日重新定价\par}
{\small\color{refgray}数据：VGT报价从5月28日高点回调51\%，同期恒生科技指数仅下跌5\%\par}
{\small\color{refgray}数据：2026-2028年盈利预测下调5\%-13\%，目标价从600港元调整至500港元\par}
{\small\color{refgray}边际变化：市场此前关注CoWoS-L路线图推进，但JPM研讨会指出5.5倍尺寸已遇困难，9.5倍焊接问题在2027年上半年仍待解决，封装密度提升节奏可能被迫调整。企业IT支出方面，内存成本上升正重塑OEM定价与供应格局，报价有效期大幅缩短推动客户提前下单。\par}
\subsubsection*{关键数据}
\begin{itemize}
  \item Kimi K3参数规模2.8T，是DeepSeek-V3（671B）的4倍以上，在MXFP4下仍需1.4TB HBM
  \item CoWoS-L光罩尺寸从3.3倍扩大到9.5倍，基板尺寸从85mm见方增至130mm x 140mm
  \item OEM报价有效期从60-90天缩短至1-2周，存储和AI集群交付周期超6个月
  \item 康宁二季度营收超预期幅度约5000万美元，太阳能业务对每股收益影响约0.07美元
  \item VGT报价从5月28日高点回调51\%，同期恒生科技指数仅下跌5\%
  \item 2026-2028年盈利预测下调5\%-13\%，目标价从600港元调整至500港元
  \item 全球token使用量自2025年以来每周平均增长6\%，自2026年以来每周增长9\%
  \item 多智能体工作负载每个查询产生的输出token比传统非推理聊天工作负载多20万倍
\end{itemize}
\begin{figure}[htbp]
\centering
\includegraphics[width=0.88\linewidth]{figures/fig_017.jpg}
\caption{Exhibit 1 - Exhibit 1: Global token usage increased on average by +6\% per week since 2025, or +9\% per week since 2026, with China models accounting for \textbackslash{}\textasciitilde{}70\% of the tokens in July 2026 Weekly Token Usage by Major Models (API) via OpenRouter}
\end{figure}
\begin{figure}[htbp]
\centering
\includegraphics[width=0.88\linewidth]{figures/fig_018.jpg}
\caption{Exhibit 3 - Exhibit 3: Open MoE model total parameters and active parameters per token continue to scale upward, resulting in larger weight for memory storage Open MoE models total parameters (bn) vs. active parameters per token (bn) BofA GL}
\end{figure}
\begin{figure}[htbp]
\centering
\includegraphics[width=0.88\linewidth]{figures/fig_019.jpg}
\caption{Exhibit 4 - Exhibit 4: Latest multi-agent workloads generate up to 200,000x more output tokens per query than traditional non-reasoning chat workloads (GPT-4o) Tokens per query by different model workloads BofA GLOBAL RESEARCH Exhibit 5: New}
\end{figure}
\begin{figure}[htbp]
\centering
\includegraphics[width=0.88\linewidth]{figures/fig_048.jpg}
\caption{Exhibit 2 - Exhibit 2: Industry Average Utilization at Display Fabs Long-Term Beneficiary of Architecture-Agnostic Scale-Up Adoption. As discussed in our refreshed Scale-Up Network primer, recent volatility across the optical ecosystem has r}
\end{figure}
{\small\color{refgray}本节综合 5 篇研究；涉及 BofA、MS、JPM。}
\newpage
\subsection{中国工业与材料：库存分化与结构性升级}
\textbf{中国工业与材料板块呈现显著的结构性分化：电池材料受益于储能高景气与锂供给约束，铜铝库存去化加速，而工程机械出货量转正但利用率仍弱，煤炭则受季节性高温驱动反弹。AI数据中心电源与高端农机成为新的增长亮点。}
\subsubsection*{主流共识}
\begin{itemize}
  \item 储能电池需求强劲，2026年上半年产能利用率达满产，全年出货目标维持高位。
  \item 铜铝库存持续去化，中国境内铜库存降至近十年同期最低，洋山铜溢价回升至100美元/吨。
  \item 动力煤价格受夏季高温推动反弹，CCI 5500指数周环比上涨2.6\%至821元/吨。
  \item AI数据中心电源设备需求周期明确，订单积压扩大，领先供应商定价能力维持。
\end{itemize}
\subsubsection*{分歧与条件差异}
\begin{itemize}
  \item 锂市场：2027年供需缺口预期与津巴布韦出口禁令执行的不确定性，分歧在于供给增量能否如期释放。
  \item 工程机械：塔吊出货量同比转正（+11\%），但机队利用率连续四个月下降，分歧在于需求复苏的可持续性。
  \item 铜市场：中国去库与美国囤货形成镜像，分歧在于驱动因素——中国是供需再平衡，美国是政策预期套利。
  \item 电池产业：韩国电池企业面临中国竞争对手在欧洲市场份额扩大的压力，分歧在于其储能转型能否弥补EV份额损失。
\end{itemize}
\subsubsection*{机构观点}
\paragraph{Citi} 2027年全球锂需求增量预计超过供给增量60-110千吨，对锂价持积极看法，但前提是津巴布韦锂精矿出口禁令如期执行。当前港口锂精矿库存周环比下降3.9万吨至32.7万吨，冶炼厂惜售，正极材料厂预计8月开始补库。塔吊出货量6月同比+11\%，为2023年7月以来首次转正，但机队利用率连续四个月下降，出货量与利用率背离。国轩高科上半年电池出货约65GWh，储能电池产能利用率100\%，EV电池约70\%，全年出货目标维持150GWh。
{\small\color{refgray}数据：2027年全球锂需求增量400-450千吨，供给增量约340千吨\par}
{\small\color{refgray}数据：津巴布韦禁令若执行，供给减少约70千吨\par}
{\small\color{refgray}数据：国内港口锂精矿库存32.7万吨，周环比-3.9万吨\par}
{\small\color{refgray}数据：塔吊出货量6月同比+11\%\par}
{\small\color{refgray}边际变化：锂市场从供给过剩叙事转向供需缺口预期，塔吊出货量转正但利用率仍弱，显示复苏不均衡。\par}
\paragraph{MS} 动力煤价格因夏季高温反弹，CCI 5500周环比+2.6\%至821元/吨，产地大同5800坑口价+3.7\%至708元/吨，但焦煤价格走弱，柳林主焦煤-1.7\%至850元/吨。WAIC 2026显示物理AI快速迭代，工业AI从监控向工程自动化扩展，人形机器人进入产线，但数据不足是主要挑战。韩国电池企业正从单一依赖EV向EV+储能双轨转型，美国ESS需求预计2030年达279GWh，但中国电池厂在欧洲份额持续扩大。
{\small\color{refgray}数据：CCI 5500指数821元/吨，周环比+2.6\%\par}
{\small\color{refgray}数据：秦皇岛港5500大卡722元/吨，微涨0.1\%\par}
{\small\color{refgray}数据：柳林4号主焦煤850元/吨，周环比-1.7\%\par}
{\small\color{refgray}数据：WAIC 2026参展企业超1000家，机器人/智能硬件占24\%\par}
{\small\color{refgray}边际变化：动力煤反弹由季节性高温驱动，而非终端需求全面回暖；物理AI从概念进入可部署阶段，拉动工业资本开支。\par}
\paragraph{GS} 第一拖拉机受益于高端拖拉机需求结构性升级，2Q26中高端马力段产量同比+22\%，高端产品占比达37\%的历史新高，出口增长40\%。应流股份两机业务产能扩张推进，2Q26产值约4.2亿元，三台新设备下半年投产，市场关注点从公司执行转向燃气轮机订单周期持续性，但替换周期和数据中心增量需求仍支撑订单。
{\small\color{refgray}数据：2Q26中高端拖拉机产量同比+22\%，高端+41\%\par}
{\small\color{refgray}数据：高端产品占比37\%，历史均值约20\%\par}
{\small\color{refgray}数据：中高端拖拉机出口同比+40\%，高端出口+56\%\par}
{\small\color{refgray}数据：应流2Q26两机产值约4.2亿元，6月单月1.5亿元\par}
{\small\color{refgray}边际变化：拖拉机行业从总量增长转向产品结构升级，高端占比弹性较高；燃气轮机订单担忧从公司执行转向周期持续性，但订单积压仍在扩大。\par}
\paragraph{JPM} 中国铜铝库存去化连续五周强劲，铜库存降至约12万吨（近十年同期最低），洋山铜溢价回升至100美元/吨。全球铜库存结构分化：中国去库45万吨，美国囤货至72.5万吨（占全球70\%）。德业股份上半年净利润同比+75\%，新兴市场贡献约60\%收入，管理层可能上调全年指引。潍柴动力AIDC业务执行符合预期，2.5MW燃气发电机已发布，美国EPA认证预计8月完成，市场对产能过剩担忧可能夸大。
{\small\color{refgray}数据：中国铜库存约12万吨，近十年同期最低\par}
{\small\color{refgray}数据：洋山铜溢价100美元/吨，2025年5月以来首次\par}
{\small\color{refgray}数据：铝库存单周去化5.4万吨，春节后最强\par}
{\small\color{refgray}数据：德业股份上半年净利润约27亿元，同比+75\%\par}
{\small\color{refgray}边际变化：铜铝库存去化加速，消费信号改善；全球铜库存结构从中国去库与美国囤货的镜像分化，转向关注政策落地后的去库节奏。\par}
\subsubsection*{关键数据}
\begin{itemize}
  \item 2027年全球锂需求增量400-450千吨，供给增量约340千吨
  \item 国内港口锂精矿库存32.7万吨，周环比-3.9万吨
  \item 塔吊出货量6月同比+11\%，为2023年7月以来首次转正
  \item 国轩高科上半年电池出货65GWh，储能产能利用率100\%
  \item CCI 5500指数821元/吨，周环比+2.6\%
  \item 2Q26高端拖拉机占比37\%，历史均值约20\%
  \item 中国铜库存约12万吨，近十年同期最低
  \item 洋山铜溢价100美元/吨，2025年5月以来首次
  \item 德业股份上半年净利润约27亿元，同比+75\%
  \item 美国数据中心2025年交付55GW，新订单100GW
\end{itemize}
\begin{figure}[htbp]
\centering
\includegraphics[width=0.88\linewidth]{figures/fig_074.jpg}
\caption{Exhibit 1 - Exhibit 2: China's medium-to-high HP tractors production volume saw strong growth of \$+22\textbackslash{}\%\$ yoy in 2Q26, with high-HP tractors recording robust rebound of \$+41\textbackslash{}\%\$ yoy China quarter medium-to-high HP tractor production volume (k un}
\end{figure}
\begin{figure}[htbp]
\centering
\includegraphics[width=0.88\linewidth]{figures/fig_087.jpg}
\caption{Figure 1 - Figure 2: Aluminum inventories movements in 2026 – weekly change in China visible aluminium inventory (SHFE + Bonded). 52kt of Aluminium destocking in the past week, strongest since 2026 CNY x-axis = weeks post Chinese New Year (w}
\end{figure}
\begin{figure}[htbp]
\centering
\includegraphics[width=0.88\linewidth]{figures/fig_033.jpg}
\caption{Exhibit 1 - Exhibit 1: Chart of the Week – Daily Thermal Coal Usage at 25 Provinces \#\# CHINA COAL}
\end{figure}
\begin{figure}[htbp]
\centering
\includegraphics[width=0.88\linewidth]{figures/fig_116.jpg}
\caption{Figure 1 - Figure 1: HDT monthly sales volume and Y/Y change \% Figure 2: LNG and Diesel price gap Figure 3: China's electric HDT domestic penetration (\%)}
\end{figure}
{\small\color{refgray}本节综合 12 篇研究；涉及 Citi、MS、GS、JPM。}
\newpage
\subsection{中国消费与房地产：分化加剧，龙头韧性凸显}
\textbf{中国消费与房地产领域呈现显著分化：房地产销售在政策支持下出现阶段性企稳，但城市层级和公司类型间差异巨大；消费领域，茶饮品牌开店节奏分化，高端白酒定价能力凸显，黄金珠宝需求结构性改善。整体来看，行业龙头凭借结构性优势在调整中保持韧性，而中小企业面临更大压力。}
\subsubsection*{主流共识}
\begin{itemize}
  \item 房地产销售出现边际改善，新房和二手房周度成交同比转正，但城市间分化明显，二线城市表现优于一线。
  \item 消费领域，加盟商信心成为茶饮品牌扩张的关键约束，头部品牌开店指引下调，行业从高速扩张转向分化。
  \item 高端白酒（如茅台）通过主动调价管理供需关系，批价价差收窄，渠道库存结构优化。
  \item 黄金珠宝需求结构性改善，按克计价产品稳健，高端定价产品出现回暖迹象，港澳市场持续活跃。
\end{itemize}
\subsubsection*{分歧与条件差异}
\begin{itemize}
  \item 房地产销售改善的可持续性存在分歧：部分机构认为政策效果逐步显现，但去化率大幅波动和一线城市新房成交转弱表明市场信心仍不稳定。
  \item 茶饮品牌中，瑞幸凭借直营模式逆势扩张，而古茗、蜜雪冰城等加盟模式品牌开店节奏放缓，市场对盈利路径看法不一。
  \item 黄金珠宝领域，周大福和老铺黄金短期趋势分化：周大福同店销售增长稳健，老铺黄金仍在调整，但市场对三季度可见度存在分歧。
\end{itemize}
\subsubsection*{机构观点}
\paragraph{Citi} 房地产行业上半年业绩分化加剧，央企利润占比升至近九成，民企仍在深度调整。19家样本企业上半年核心净利润合计约178亿元，同比降32\%，但央企仅降20\%，地方国企降82\%，民企由盈转亏。销售端已有回暖信号，5家公司实现6\%-15\%同比增长，主要受益于核心城市复苏。三季度是关键观察期，预计更多房企销售增速转正。
{\small\color{refgray}数据：19家样本企业1H26核心净利润合计约178亿元，同比降32\%\par}
{\small\color{refgray}数据：6家中央国企核心净利润约187亿元，占样本总利润105\%\par}
{\small\color{refgray}数据：5家公司实现合约销售正增长：中海+12\%、金茂+8\%、招商蛇口+8\%、华润置地+6\%、中国金茂+15\%\par}
{\small\color{refgray}数据：前5月重点房企拿地金额同比-57\%，300城土地供应创20年新低\par}
{\small\color{refgray}边际变化：相比此前市场对上半年业绩的悲观预期，实际盈利降幅好于预期，且销售端在二季度出现边际改善，央企韧性超出预期。\par}
\paragraph{MS} 黄金珠宝行业6月数据呈现分化：按克计价产品同店销售保持健康但环比走弱，溢价固定价黄金需求同比改善。六福中国2Q同店销售增长16\%，低于4-5月超20\%增速；港澳市场强劲，2Q同店销售增长41\%。周大福1QFY27同店销售趋势与4-5月基本一致（中国+20\%，港澳+41\%），金价稳定和人民币升值提供支撑。老铺黄金仍在调整，但2Q趋势已大部分被市场消化。
{\small\color{refgray}数据：六福中国2Q同店销售增长16\%，低于4-5月超20\%增速\par}
{\small\color{refgray}数据：六福港澳2Q同店销售增长41\%，与4-5月持平\par}
{\small\color{refgray}数据：周大福1QFY27中国同店销售增长约20\%，港澳约41\%\par}
{\small\color{refgray}数据：当前金价约4000美元，低于公司假设的4300美元\par}
{\small\color{refgray}边际变化：6月数据显示按克计价产品环比走弱，但高端定价产品出现回暖迹象，市场对老铺黄金的调整预期已部分反映在定价中。\par}
\paragraph{MS} 截至7月19日当周，50城新房周度成交同比+3\%，扭转前一周+1\%的微弱增幅；10城二手房成交同比+6\%，较前一周-3\%明显改善。新房成交转正主要由二线城市拉动（同比+14\%），但一线城市同比-24\%，与前一周+40\%形成反差。二手房市场修复更稳定，一线城市同比+9\%，二线城市同比+5\%。年初至今10城二手房累计同比+4\%，50城新房累计仍为-11\%。
{\small\color{refgray}数据：50城新房周度成交同比+3\%（前一周+1\%）\par}
{\small\color{refgray}数据：10城二手房周度成交同比+6\%（前一周-3\%）\par}
{\small\color{refgray}数据：二线城市新房成交同比+14\%（前一周0\%）\par}
{\small\color{refgray}数据：一线城市新房成交同比-24\%（前一周+40\%）\par}
{\small\color{refgray}边际变化：新房和二手房周度成交均出现同比正增长，但一线城市新房成交大幅转弱，去化率从95\%骤降至26\%（受推盘区位影响），市场信心仍不稳定。\par}
\paragraph{JPM} 茶饮品牌开店节奏分化加剧，加盟商信心成为扩张关键约束。古茗FY26净开店目标下调43\%，蜜雪扩张放缓更明显，盈利预测下调11\%。瑞幸逆势扩张，预计1H26净开4700家店，直营模式（2.2万家直营占总门店三分之二）提供运营控制力。行业从'水涨船高'转向'分屏格局'，加盟商对客流变化的关注上升，下半年盈利存在节奏变化不确定性。
{\small\color{refgray}数据：古茗FY26净开店目标下调43\%\par}
{\small\color{refgray}数据：蜜雪冰城FY26盈利预测下调11\%\par}
{\small\color{refgray}数据：瑞幸预计1H26净开4700家店\par}
{\small\color{refgray}数据：瑞幸直营店约2.2万家，占总门店三分之二\par}
{\small\color{refgray}边际变化：相比此前市场对茶饮行业高速扩张的预期，加盟商信心拐点已确认，古茗和蜜雪开店指引大幅下调，而瑞幸凭借直营模式保持逆势扩张。\par}
\paragraph{JPM} 茅台2026年7月宣布年内第二次调价，出厂价从1269元升至1369元（+7.9\%），直销零售价从1539元升至1639元（+6.5\%），涨幅对称。调价节奏加快（距3月仅隔不到四个月）且选择淡季执行，显示公司对终端需求信心强于市场预期。批价数据显示整箱与散瓶价差显著收窄，投机性囤货减少，真实消费占比上升。老酒库存价值同步重估，但囤货风险增加，推动库存向真实消费流转。
{\small\color{refgray}数据：出厂价从1269元升至1369元，涨幅7.9\%\par}
{\small\color{refgray}数据：直销零售价从1539元升至1639元，涨幅6.5\%\par}
{\small\color{refgray}数据：2026年内第二次调价，距3月仅隔不到四个月\par}
{\small\color{refgray}数据：整箱与散瓶批价价差显著收窄\par}
{\small\color{refgray}边际变化：相比3月不对称调价（出厂价涨幅大于直销价），此次对称调价显示渠道利润分配格局稳定，且调价节奏加快表明公司主动管理供需关系，从被动跟随市场价转向主动锚定。\par}
\subsubsection*{关键数据}
\begin{itemize}
  \item 19家样本房企1H26核心净利润合计约178亿元，同比降32\%，央企利润占比升至105\%
  \item 50城新房周度成交同比+3\%，10城二手房同比+6\%，年初至今二手房累计+4\%
  \item 古茗FY26净开店目标下调43\%，瑞幸预计1H26净开4700家店
  \item 茅台出厂价从1269元升至1369元，2026年内第二次调价
  \item 六福中国2Q同店销售增长16\%，港澳2Q增长41\%
  \item 前5月重点房企拿地金额同比-57\%，300城土地供应创20年新低
\end{itemize}
\begin{figure}[htbp]
\centering
\includegraphics[width=0.88\linewidth]{figures/fig_024.jpg}
\caption{Figure 2 - Figure 2. Listed Names 5M26 Land Acq Attri. Cost -57\%yoy © 2026 Citi Inc. No redistribution without Citi's written permission. Figure 3. Listed Names 1H26 Contracted Sales -15\% yoy © 2026 Citi Inc. No redistribution without Citi's written permission. \#\# Sector}
\end{figure}
\begin{figure}[htbp]
\centering
\includegraphics[width=0.88\linewidth]{figures/fig_037.jpg}
\caption{Exhibit 1 - Exhibit 1: Primary weekly unit sales and four-week moving average – Total 50 cities MS ASIA LIMITED+ Stephen Cheung, CFA}
\end{figure}
\begin{figure}[htbp]
\centering
\includegraphics[width=0.88\linewidth]{figures/fig_092.jpg}
\caption{Figure 2 - Figure 2: 2026 YTD store net openings comparison - Luckin ahead of the pack Figure 3: Store count vs average ticket size – tea brands Note: As of end-Jun 2026. Figure 4: Store count vs average ticket size – coffee brands Note:}
\end{figure}
\begin{figure}[htbp]
\centering
\includegraphics[width=0.88\linewidth]{figures/fig_107.jpg}
\caption{Figure 1 - Figure 1: Price gap of Feitian (case-packed) and single-bottle significantly narrowed after recent price hikes. Distribution price (Rmb/500ml)}
\end{figure}
{\small\color{refgray}本节综合 5 篇研究；涉及 Citi、MS、JPM。}
\newpage
\subsection{医疗保健：仿制药板块盈利稳健，创新药管线催化密集}
\textbf{仿制药板块受益于基本面改善与资金轮动，估值仍有吸引力；创新药领域，DLL3 ADC与RAS抑制剂数据更新指向差异化机会，但样本量小、随访期短仍是主要不确定性。}
\subsubsection*{主流共识}
\begin{itemize}
  \item 仿制药板块近几个季度跑赢医疗板块和整体市场，核心驱动力来自盈利执行、管线推进和宏观资金轮动。
  \item TEVA、VTRS等仿制药企二季度业绩预计符合预期，市场更关注下半年催化剂事件（如PDUFA、Ph3数据）。
  \item DLL3 ADC在SCLC一线治疗中，PFS和安全性改善比ORR更重要，去化疗双药方案是差异化方向。
  \item RAS抑制剂ERAS-0015在胰腺癌中展现出剂量-应答关系，安全性支持长期用药，但数据仍基于小样本。
\end{itemize}
\subsubsection*{分歧与条件差异}
\begin{itemize}
  \item TEVA二季度每股收益预测分歧：GS预计0.09美元，显著低于市场共识0.16美元，差异来自研发费用（Emalex收购未完全反映）。
  \item Zoci在ESMO的数据解读分歧：市场可能更关注ORR，但Citi认为PFS延长至“高个位数月”才是真正突破。
  \item ERAS-0015的注册路径：MS认为需更成熟数据（2026H2-2027H1验证），而公司已提前部分适应症上市时间。
\end{itemize}
\subsubsection*{机构观点}
\paragraph{Bernstein} 仿制药板块二季度业绩预计整体符合预期，但TEVA每股收益预测（0.09美元）低于市场共识（0.16美元），主要因研发费用较高（Emalex收购尚未完全反映）。市场焦点转向下半年催化剂：长效奥氮平（10月PDUFA）、Ecopipam（FDA申报）、DARI（年底Ph3数据）。VTRS受益于商业业务改善和估值折价，关注Nashik工厂产能恢复及中国业务表现。板块估值仍有吸引力，但宏观轮动可能带来波动。
{\small\color{refgray}数据：TEVA二季度收入预测40.07亿美元（市场共识40.43亿美元）\par}
{\small\color{refgray}数据：TEVA调整后每股收益预测0.09美元（市场预期0.16美元）\par}
{\small\color{refgray}数据：VTRS年初至今上涨约30\%（vs标普500）\par}
{\small\color{refgray}数据：TEVA年初至今下跌约8\%\par}
{\small\color{refgray}边际变化：从关注二季度业绩转向下半年催化剂密集期，管线进展成为估值重估核心。\par}
\paragraph{Citi} 再鼎医药Zoci（DLL3 ADC）将在ESMO首次公布一线SCLC数据，预计入组30-40例患者。核心看点不是ORR（现有标准疗法已达60-68\%），而是PFS能否实现“高个位数月”改善及安全性优势。去化疗双药方案（Zoci+阿替利珠单抗）因耐受性更好、治疗持续时间更长，可能成为差异化方向。CNS活性（脑转移控制）是额外区分点。
{\small\color{refgray}数据：现有标准疗法ORR约60-68\%，PFS约5个月\par}
{\small\color{refgray}数据：3级以上不良事件约62-67\%\par}
{\small\color{refgray}数据：预计入组30-40例一线SCLC患者\par}
{\small\color{refgray}数据：PFS改善目标：高个位数月（非仅1-2个月）\par}
{\small\color{refgray}边际变化：从二线数据转向一线治疗，去化疗双药方案成为市场关注焦点。\par}
\paragraph{MS} Erasca的ERAS-0015在2L+ KRAS G12X胰腺癌中，32mg剂量组uORR从50\%提升至57\%（4/7），24-32mg合并组从27\%升至40\%（6/15）。剂量-应答关系清晰，安全性支持长期用药（中位相对剂量强度100\%），未出现因副作用停药。但样本量小、随访期短，关键验证节点在2026H2-2027H1。联合用药（EGFR/PD-1/PRMT5）是差异化方向。
{\small\color{refgray}数据：32mg剂量组uORR 57\%（4/7）\par}
{\small\color{refgray}数据：24-32mg合并组uORR 40\%（6/15）\par}
{\small\color{refgray}数据：72\%患者出现皮疹（3级仅3\%）\par}
{\small\color{refgray}数据：中位相对剂量强度100\%\par}
{\small\color{refgray}边际变化：从早期概念验证转向多适应症注册潜力，但需更成熟数据确认疗效持久性。\par}
\subsubsection*{关键数据}
\begin{itemize}
  \item TEVA二季度每股收益预测0.09美元（市场预期0.16美元）
  \item VTRS年初至今上涨30\%
  \item Zoci一线SCLC预计入组30-40例
  \item ERAS-0015 32mg组uORR 57\%（4/7）
  \item 现有SCLC标准疗法ORR 60-68\%，PFS约5个月
  \item ERAS-0015中位相对剂量强度100\%
  \item TEVA长效奥氮平PDUFA日期10月
  \item ERAS-0015 24-32mg合并组uORR 40\%（6/15）
\end{itemize}
\begin{figure}[htbp]
\centering
\includegraphics[width=0.88\linewidth]{figures/fig_007.jpg}
\caption{Exhibit 1 - Exhibit 1: GSe vs. consensus All values in \textbackslash{}\$mn Exhibit 2: Austedo scripts Continued momentum Exhibit 3: Uzedy scripts Growing momentum}
\end{figure}
\begin{figure}[htbp]
\centering
\includegraphics[width=0.88\linewidth]{figures/fig_043.jpg}
\caption{Exhibit 1 - Exhibit 1: Efficacy data for 2L+ KRAS G12X PDAC pts across April and May data cutoffs Exhibit 2: Updated time on treatment for 2L+ KRAS G12X PDAC pts Exhibit 3: ERAS-0015 efficacy vs. comparator pan-RAS molecular glue RMC-6236}
\end{figure}
\begin{figure}[htbp]
\centering
\includegraphics[width=0.88\linewidth]{figures/fig_008.jpg}
\caption{Exhibit 2 - Exhibit 2: Austedo scripts Continued momentum Exhibit 3: Uzedy scripts Growing momentum VTRS (Neutral): Expect in-line quarter ahead of key pipeline updates in YE26+ VTRS shares have outperformed YTD (+30\% vs. S\&P 500; +36\% v}
\end{figure}
\begin{figure}[htbp]
\centering
\includegraphics[width=0.88\linewidth]{figures/fig_045.jpg}
\caption{Exhibit 2 - Exhibit 2: Updated time on treatment for 2L+ KRAS G12X PDAC pts Exhibit 3: ERAS-0015 efficacy vs. comparator pan-RAS molecular glue RMC-6236 benchmark Exhibit 4: ERAS-0015 safety vs. comparator pan-RAS molecular glue RMC-6236 b}
\end{figure}
{\small\color{refgray}本节综合 3 篇研究；涉及 Bernstein、Citi、MS。}
\newpage
\subsection{全球汽车市场分化：日本车企Q1业绩分化，中国新能源出口创新高}
\textbf{全球汽车市场呈现结构性分化：日本车企在2026财年Q1销量表现分化，本田、铃木受益于区域需求增长和产品周期，而丰田、日产面临压力；中国新能源车市场内需疲软但出口强劲，新车型订单驱动短期增长，但折扣收窄显示价格竞争降温。}
\subsubsection*{主流共识}
\begin{itemize}
  \item 日本车企Q1销量整体同比增长3.9\%，但内部分化明显，本田和铃木表现突出，丰田、日产、马自达下滑。
  \item 中国新能源车市场内需疲软，6月零售同比下降23\%，但出口创历史新高，同比增长80\%。
  \item 新能源车渗透率高位小幅回落，6月保险渗透率60.3\%，低于5月的62.3\%，属于季节性波动。
  \item 中国新能源车终端折扣收窄，经销商平均折扣率从7.29\%降至6.94\%，反映产品周期切换而非需求全面回暖。
\end{itemize}
\subsubsection*{分歧与条件差异}
\begin{itemize}
  \item 日本车企利润前景：Bernstein认为本田、铃木、丰田通商Q1利润有望超预期，而日产和马自达利润指引完成度较低，存在不确定性。
  \item 中国新能源车增长驱动力：GS强调新车型订单（小鹏Mona L03、理想L6改款）驱动短期增长，而JEF认为出口是上半年主要增长引擎，内需结构性问题更突出。
  \item 丰田销量展望：Bernstein预计丰田Q1销量下滑，但7-9月生产计划显示产量同比增长3.4\%，销量动能有望恢复；市场对丰田下半年复苏节奏存在分歧。
\end{itemize}
\subsubsection*{机构观点}
\paragraph{Bernstein} 日本车企2026财年Q1（4-6月）合计销量预计同比增长3.9\%，但内部分化显著。铃木以18.4\%的同比增速领跑，受益于印度GST下调后需求释放及Kharkhoda新工厂投产；本田全球销量增长5.9\%，美国HEV（CR-V）创纪录，日本税改后注册车需求回升，关税成本从27.5\%降至15\%改善利润空间。丰田通商非洲陆地巡洋舰销量预计增长28.4\%，利润可能超预期。丰田、日产、马自达销量同比下滑，但丰田7-9月生产计划显示产量将增长3.4\%，动能有望恢复。行业Q1营业利润预计同比增长6.5\%。
{\small\color{refgray}数据：铃木全球销量预计增长18.4\%，印度销量增长33\%\par}
{\small\color{refgray}数据：本田全球销量增长5.9\%，美国HEV销量增长8.4\%\par}
{\small\color{refgray}数据：丰田通商非洲陆地巡洋舰销量预计增长28.4\%\par}
{\small\color{refgray}数据：行业Q1营业利润预计同比增长6.5\%\par}
{\small\color{refgray}边际变化：相比此前对日本车企利润的谨慎预期，Bernstein上调本田、铃木、丰田通商Q1利润预期，认为关税成本下降和区域需求改善是超预期主因。\par}
\paragraph{GS} 2026年7月第三周，中国主要新能源车企合计订单环比增长31\%，同比增长5\%，由新车型集中交付驱动。小鹏Mona L03单周录得5.92万辆不可退订单，70\%为纯电版；理想L6改款上市带动理想订单环比增长131\%；小米SU7订单环比增长32\%。7月1-12日乘用车零售同比下降15\%，较上半年20\%降幅收窄；新能源零售渗透率63.1\%，批发渗透率69.1\%。经销商平均折扣率从7.29\%降至6.94\%，比亚迪王朝系列折扣率从4.10\%降至3.95\%，价格竞争短期降温。
{\small\color{refgray}数据：小鹏Mona L03单周5.92万辆不可退订单，纯电版占70\%\par}
{\small\color{refgray}数据：理想L6改款带动理想订单环比增长131\%\par}
{\small\color{refgray}数据：7月1-12日乘用车零售同比下降15\%\par}
{\small\color{refgray}数据：新能源零售渗透率63.1\%，批发渗透率69.1\%\par}
{\small\color{refgray}边际变化：相比此前新能源车订单增速放缓的预期，GS指出新车型（小鹏Mona L03、理想L6改款）集中交付驱动订单环比大幅增长，但折扣收窄显示价格竞争烈度下降，需求改善更多来自产品周期而非整体回暖。\par}
\paragraph{JEF} 中国电动汽车市场呈现出口强劲、内需疲软的二元格局。6月乘用车出口90.45万辆，同比增长80\%，环比增长12\%，创历史新高；国内零售仅166万辆，同比下滑23\%。1-6月出口累计增长72\%，零售下降24\%。新能源保险渗透率60.3\%，低于5月高点62.3\%，纯电41.6\%，插混18.7\%。出口由奇瑞、比亚迪、吉利、上汽主导，吉利出口份额提升5.8个百分点。中国品牌在欧洲乘用车份额9.1\%，新能源占比19.7\%；拉美份额16.3\%，新能源占比87\%。
{\small\color{refgray}数据：6月乘用车出口90.45万辆，同比增长80\%\par}
{\small\color{refgray}数据：6月国内零售166万辆，同比下滑23\%\par}
{\small\color{refgray}数据：1-6月出口累计增长72\%，零售下降24\%\par}
{\small\color{refgray}数据：新能源保险渗透率60.3\%，低于5月62.3\%\par}
{\small\color{refgray}边际变化：相比此前市场关注内需渗透率攀升，JEF强调出口成为上半年主要增长引擎，内需结构性问题（零售持续下滑）比渗透率波动更值得关注，且中国品牌在海外市场（欧洲、拉美）份额持续扩张。\par}
\subsubsection*{关键数据}
\begin{itemize}
  \item 铃木全球销量预计增长18.4\%，印度销量增长33\%
  \item 本田全球销量增长5.9\%，美国HEV销量增长8.4\%
  \item 小鹏Mona L03单周5.92万辆不可退订单，纯电版占70\%
  \item 7月1-12日新能源零售渗透率63.1\%，批发渗透率69.1\%
  \item 6月乘用车出口90.45万辆，同比增长80\%
  \item 6月国内零售166万辆，同比下滑23\%
  \item 新能源经销商平均折扣率从7.29\%降至6.94\%
  \item 吉利出口份额较去年同期提升5.8个百分点
\end{itemize}
\begin{figure}[htbp]
\centering
\includegraphics[width=0.88\linewidth]{figures/fig_012.jpg}
\caption{Exhibit 4 - EXHIBIT 2: Suzuki's estimated FY3/27 Q1 sales volume reached 893k units (+18.4\% YoY) Note: Apr-May 2026 figures are based on officially disclosed company data, while Jun figures are based on publicly available country-level data}
\end{figure}
\begin{figure}[htbp]
\centering
\includegraphics[width=0.88\linewidth]{figures/fig_069.jpg}
\caption{Exhibit 2 - Exhibit 2: NEV dealer discount trend By brand dealer discount is based on best-selling models. Exhibit 3: ICE dealer discount trend}
\end{figure}
\begin{figure}[htbp]
\centering
\includegraphics[width=0.88\linewidth]{figures/fig_111.jpg}
\caption{Exhibit 1 - Exhibit 1 - China NEV Monthly Insurance Sales Penetration Reached 60.3\% in Jun-26 Exhibit 2 - Chery and BYD Led China's 1H26 PV Exports, Leapmotor Posted Sharpest YoY Gain Exhibit 3 - 1H26 Exports (+72\%) and NEVs (+7\%) Offset We}
\end{figure}
\begin{figure}[htbp]
\centering
\includegraphics[width=0.88\linewidth]{figures/fig_113.jpg}
\caption{Exhibit 6 - Exhibit 6 - PV Exports Improved 80\% YoY to 904.5k Units Exhibit 7 - NEV Exports Improved 113\% YoY to 434.7k Units Exhibit 8 - Export Volume by Major Domestic Brand}
\end{figure}
{\small\color{refgray}本节综合 3 篇研究；涉及 Bernstein、GS、JEF。}
\newpage
\subsection{其他行业与主题：结构性分化与政策驱动下的投资机会}
\textbf{本期其他行业与主题覆盖清洁电器、物业、造船、天然气、日本国防与财政、韩国银行及全球权益扩散等多个领域，核心线索在于：市场正从对宏观风险的过度担忧转向对结构性变化的重新定价，包括中国清洁电器国内需求超预期、日本财政框架转向扩张、欧洲天然气供应收紧推升价格、以及全球权益从集中走向扩散。}
\subsubsection*{主流共识}
\begin{itemize}
  \item 中国清洁电器618销售增速环比改善，国内需求好于市场担忧，海外增长保持韧性。
  \item 日本财政框架从单年度预算转向跨年度预算，为中长期增长战略创造资金空间。
  \item 欧洲天然气供应因霍尔木兹海峡出口延迟而收紧，冬季储气库缓冲空间被压缩。
  \item 全球权益市场正从少数头部公司向更广泛的板块扩散，动量因子出现极端反转。
\end{itemize}
\subsubsection*{分歧与条件差异}
\begin{itemize}
  \item 清洁电器利润率前景：GS认为成本控制与定价稳定可对冲压力，但市场仍担忧追觅竞争与汇率影响。
  \item 日本财政可持续性：GS模拟显示低利率下财政扩张有十年缓冲期，但若利率升至3.5\%，债务比率将快速恶化。
  \item 欧洲天然气价格路径：GS上调TTF预测至60/53欧元/兆瓦时，但若出口恢复快于预期，价格可能回落至40欧元。
  \item 中国科技能力提升：JEF认为是结构性转变，但市场对其对西方企业的竞争影响存在分歧。
\end{itemize}
\subsubsection*{机构观点}
\paragraph{GS} 清洁电器板块年初至今下跌约30\%，市场担忧过度。618期间石头科技和科沃斯GMV同比增长超20\%，石头科技在京东/天猫扫地机器人市占率同比提升10个百分点至35\%。海外亚马逊美国站销售额同比增长39\%，石头科技市占率从20\%升至32\%。利润率方面，国内扫地机器人均价坚挺，两家公司已停止补贴消费者，成本与汇率压力部分被关税退税和亏损业务调整对冲。
{\small\color{refgray}数据：石头科技618扫地机器人市占率35\%，同比+10ppt\par}
{\small\color{refgray}数据：亚马逊美国站清洁电器销售额2Q26同比+39\%\par}
{\small\color{refgray}数据：石头科技亚马逊美国市占率从20\%升至32\%\par}
{\small\color{refgray}边际变化：市场此前担忧追觅激进扩张，但报告指出追觅已调整策略，行业价格竞争烈度可能低于预期。\par}
\paragraph{GS} 招商积余面临结构性矛盾：资产负债表上投资性房地产占总资产27\%，出租率同比下降2ppt，租金收入同比下降4\%，而同期整体营收增长12\%。GS预计公允价值损失将从2025财年的1900万元显著扩大。股东回报方面，当前派息率约30\%无提升计划，未来三年预期股息率仅约3\%，远低于同业6-7\%的平均水平。
{\small\color{refgray}数据：投资性房地产55亿元，占总资产27\%\par}
{\small\color{refgray}数据：出租率同比-2ppt，租金收入同比-4\%\par}
{\small\color{refgray}数据：2025财年公允价值损失1900万元\par}
{\small\color{refgray}数据：预期股息率约3\%，同业平均6-7\%\par}
{\small\color{refgray}边际变化：GS将评级从中性下调至偏谨慎，核心逻辑从短期业绩波动转向股东回报优先级偏低与资产减值不确定性叠加。\par}
\paragraph{GS} 松发股份（恒力重工）定增定价每股146元，较基准价折价约2\%，较收盘价折价约6\%，反映参与方对中长期前景乐观。三期产能已于2026年6月投产，新增约100万CGT产能，超过前两期总和。满产状态下，现有三期产能2027-2028年可贡献140-160亿元盈利。恒力重工有望2027年成为全球第二大造船商，2025-2027年产能CAGR达29\%，远超全球同行3\%。
{\small\color{refgray}数据：定增价146元，较基准价折价2\%\par}
{\small\color{refgray}数据：三期产能新增约100万CGT\par}
{\small\color{refgray}数据：满产可贡献140-160亿元盈利\par}
{\small\color{refgray}数据：2025-2027年产能CAGR 29\%\par}
{\small\color{refgray}边际变化：定增折价幅度小且三期已投产，市场对稀释影响已基本消化，关注点转向产能爬坡与全球份额提升。\par}
\paragraph{GS} 霍尔木兹海峡LNG出口正常化时间从2026年7月推迟至10月，全球供应净减少约1600万吨/年（全球4\%）。欧洲冬季储气库填充率将从预期74\%降至67\%，冬季结束时降至28\%。GS将3Q26/4Q26 TTF预测从41/40欧元上调至60/53欧元，接近远期合约58/57欧元。若出口推迟至2027年，TTF可能升至100欧元以上。
{\small\color{refgray}数据：LNG供应净减少约1600万吨/年\par}
{\small\color{refgray}数据：欧洲冬季储气库填充率降至67\%\par}
{\small\color{refgray}数据：3Q26 TTF预测上调至60欧元/兆瓦时\par}
{\small\color{refgray}数据：若出口推迟至2027年，TTF可能升至100欧元以上\par}
{\small\color{refgray}边际变化：此前预期7月恢复，现推迟至10月，导致欧洲夏季末供应缺口扩大，冬季脆弱性显著上升。\par}
\paragraph{GS} 日本骨太方针2026确立17个战略领域，从单年度预算转向跨年度预算，放弃基础财政收支盈余目标，改为通过经济增长降低债务/GDP比率。GS模拟显示，若10年期国债收益率维持在2.7\%，债务比率有约10年缓冲期；若升至3.5\%，约5年内转为上升。安保三文件修订预计2026年底前完成，将决定下一期防卫力整备计划规模与财源。
{\small\color{refgray}数据：17个战略领域，62项关键技术\par}
{\small\color{refgray}数据：目标2040年官民累计投资超370万亿日元\par}
{\small\color{refgray}数据：若10年期国债收益率2.7\%，债务比率有10年缓冲期\par}
{\small\color{refgray}数据：安保三文件修订预计2026年底前完成\par}
{\small\color{refgray}边际变化：财政框架从单年度平衡转向多年度增长导向，为国防和科技投资提供更大空间，但市场对财政纪律的担忧可能推高长期利率。\par}
\paragraph{GS} 韩国银行股被视为非科技领域的防御性选择，受益于宏观顺风（加息周期、资产质量韧性）和股东回报承诺（Value Up计划，TSR上限提高至50\%以上）。GS认为主要银行有望达到1倍市净率。保险板块受监管改革推动，但市场仍在等待更实质性的Value Up方案。金融科技板块热度有限，因政府收紧零售贷款。
{\small\color{refgray}数据：Value Up计划TSR上限提高至50\%以上\par}
{\small\color{refgray}数据：主要银行有望达到1倍市净率\par}
{\small\color{refgray}数据：保险板块受GA佣金上限和理赔审查趋严推动\par}
{\small\color{refgray}边际变化：科技股回调后，资金重新评估银行股的防御属性，叠加加息周期和股东回报升级，银行股进入有利位置。\par}
\paragraph{HSBC} 全球动量因子三周内下跌15\%，此前经历了有记录以来最强上涨之一，极端走势预示更剧烈反转。企业盈利正从科技/能源向外扩散：美国标普500盈利修正比达73\%，为2021年以来最高；新兴市场（除台韩）盈利增速预计10.5\%，为2022年以来最快；欧洲盈利增速15\%。等权重指数上涨（美国+12\%，新兴市场+5\%，欧洲+10\%）与集中度高位并存，暗示结构性轮动。
{\small\color{refgray}数据：动量因子三周内下跌15\%\par}
{\small\color{refgray}数据：美国盈利修正比73\%，为2021年以来最高\par}
{\small\color{refgray}数据：新兴市场（除台韩）盈利增速10.5\%\par}
{\small\color{refgray}数据：欧洲盈利增速15\%\par}
{\small\color{refgray}边际变化：市场从少数头部公司向更广泛板块扩散，五个C因素（企业盈利、央行、资本开支、消费者、资本流动）支撑轮动。\par}
\paragraph{JEF} 中国科技能力提升是结构性转变，非周期性波动。自然指数2026显示全球前十高影响力研究机构中九家来自中国；高被引论文全球份额从2005年约5\%升至2025年约47\%，同期美国从约40\%降至约9\%。研发支出按购买力平价已超美国（中国1.03万亿美元 vs 美国1.01万亿），年增速12\%是美国3倍多。专利申请180万件，是美国的3.6倍。DeepSeek token使用量已超过美国所有前沿模型。
{\small\color{refgray}数据：全球前十研究机构中九家来自中国\par}
{\small\color{refgray}数据：高被引论文份额中国47\% vs 美国9\%\par}
{\small\color{refgray}数据：研发支出中国1.03万亿美元 vs 美国1.01万亿\par}
{\small\color{refgray}数据：专利申请180万件，是美国的3.6倍\par}
{\small\color{refgray}边际变化：从研究产出到产业应用的多维度数据均显示中国科技实力加速提升，市场需重新评估其对全球竞争格局的影响。\par}
\paragraph{Bernstein} 地平线机器人1H26收入中点约20亿元，同比增长30\%，但较Bernstein预测低10.6\%，较彭博一致预期低15.1\%。收入低于预期主要反映中国乘用车销量变化及新设计项目获取进度慢于预期。毛利率中点63.1\%，在竞争环境下仍维持60\%以上，显示软硬件一体化定价能力。Non-GAAP净亏损14-17亿元，基本符合预期，费用管控有效。与比亚迪合作GodEye B车型即将落地，管理层预计在GodEye C方案中份额从2025年约40\%提升至2027年90\%以上。
{\small\color{refgray}数据：1H26收入中点20亿元，同比+30\%\par}
{\small\color{refgray}数据：毛利率中点63.1\%\par}
{\small\color{refgray}数据：Non-GAAP净亏损14-17亿元\par}
{\small\color{refgray}数据：GodEye C方案份额从40\%升至90\%以上\par}
{\small\color{refgray}边际变化：收入低于预期但亏损未同步扩大，费用纪律值得关注；比亚迪合作进展缓解了自研芯片替代担忧。\par}
\paragraph{MS} 百度申请将香港上市地位从第二上市转换为双重主要上市，若在2026年9月3日前完成，最早9月7日可纳入港股通。参照同类科技公司，南向资金在纳入后6-12个月可能带来30-68亿美元增量，对应7-15\%持股比例。但核心广告业务复苏有限，下半年AI投入加大，核心经营利润面临不确定性。短期事件驱动有支撑，中长期仍需收入/盈利企稳。
{\small\color{refgray}数据：南向资金增量30-68亿美元\par}
{\small\color{refgray}数据：南向持股比例7-15\%\par}
{\small\color{refgray}数据：转换流程可能仅需1-1.5个月\par}
{\small\color{refgray}边际变化：双重主要上市打开南向资金通道，成为短期催化剂，但基本面改善仍需时间。\par}
\paragraph{MS} 北京首都机场1H26净利润预计1000-2000万元，2Q26单季净利润100-1100万元，同比扭亏（2Q25净亏损3900万元）。扭亏归因于非国内客流增长，尽管总客运量环比下降4\%、同比持平。全球能源冲击推高航空票价，国内出行需求偏软。MS认为全年维持微利，2026年EPS预测仅0.01元人民币，预期PE高达78倍。
{\small\color{refgray}数据：1H26净利润1000-2000万元\par}
{\small\color{refgray}数据：2Q26单季净利润100-1100万元\par}
{\small\color{refgray}数据：总客运量环比-4\%\par}
{\small\color{refgray}数据：2026年EPS预测0.01元\par}
{\small\color{refgray}边际变化：连续两个季度盈利，但利润极薄，盈利改善路径缓慢，国际线恢复是核心变量。\par}
\subsubsection*{关键数据}
\begin{itemize}
  \item 石头科技618扫地机器人市占率35\%，同比+10ppt
  \item 招商积余投资性房地产占总资产27\%，出租率同比-2ppt
  \item 恒力重工三期产能新增约100万CGT，满产可贡献140-160亿元盈利
  \item 欧洲LNG供应净减少约1600万吨/年，3Q26 TTF预测上调至60欧元/兆瓦时
  \item 日本财政框架转向多年度预算，若10年期国债收益率2.7\%有10年缓冲期
  \item 韩国银行Value Up计划TSR上限提高至50\%以上
  \item 全球动量因子三周内下跌15\%，美国盈利修正比73\%
  \item 中国高被引论文全球份额47\% vs 美国9\%，研发支出已超美国
\end{itemize}
\begin{figure}[htbp]
\centering
\includegraphics[width=0.88\linewidth]{figures/fig_057.jpg}
\caption{Exhibit 1 - Exhibit 1: RVC/wet dry vacuum industry sales growth showed qoq improvements helped by 618. Retail sales growth by channel, yoy \% Exhibit 2: On pricing, wet dry vacuums showed more ASP pressure especially online, while RVC ASP h}
\end{figure}
\begin{figure}[htbp]
\centering
\includegraphics[width=0.88\linewidth]{figures/fig_078.jpg}
\caption{Exhibit 1 - Exhibit 1: Visible LNG exports via the Strait of Hormuz have halted... Exhibit 2: ...leading us to assume a slower ramp of Qatari LNG exports Exhibit 3: With Asia LNG demand back up yoy, at least during peak summer... Global LN}
\end{figure}
\begin{figure}[htbp]
\centering
\includegraphics[width=0.88\linewidth]{figures/fig_083.jpg}
\caption{Exhibit 1 - Exhibit 2: Higher Market Interest Rates Would Have a Large Long-Term Impact on the Public Debt-to-GDP Ratio Debt/GDP Ratio Simulation by Interest Rate Scenario, Assuming an Increase in Defense Spending and Consumption Tax Cut on Fo}
\end{figure}
\begin{figure}[htbp]
\centering
\includegraphics[width=0.88\linewidth]{figures/fig_032.jpg}
\caption{Exhibit 1 - Exhibit 1: Southbound ownership of major HSTECH constituents 30\% Exhibit 2: Primary conversion timeline for historical precedents (sorted by time to complete primary conversion)}
\end{figure}
{\small\color{refgray}本节综合 12 篇研究；涉及 Bernstein、MS、GS、HSBC、JEF。}
\newpage
\section{辅助信号｜战略咨询}
波士顿咨询2026年全球数字政府公民调查显示，公民对AI政务的接受度快速提升，但政府数字服务满意度反而下降，形成“采用率升、满意度降”的悖论。AI熟练度与正面态度强相关，公民对代理型AI的接受度超出预期，但偏好集中于外部服务而非内部决策。当前公民对政府AI部署的开放态度窗口可能不会持续，政府需抓住时机。
\subsection{AI熟练度飙升重塑公民态度，政府面临窗口期}
\textbf{公民AI熟练度两年内跃升26个百分点，熟练用户对AI正面看法是零经验用户的七倍，但政府数字服务满意度下降13个百分点，表明公民预期被私营部门拉升后形成落差。政府需在舆论固化前加速AI部署，利用熟练度杠杆塑造民意基础。}
\begin{itemize}
  \item 2026年定期使用AI工具的公民比例达64\%，较2024年上升26个百分点，对AI不确定性的“不知道”回应下降4个百分点，标志认知从模糊走向具体。
  \item 自评为“专家”的用户认为AI收益大于不确定性的比例是零经验用户的近七倍，表明真实接触是态度转变的关键，政府应优先在低不确定性场景推广AI使用。
  \item 过去十年政府数字服务使用率提升15个百分点，但公民满意度下降13个百分点，无故障使用率降5\%，矛盾源于服务复杂化、bug增多及私企拉高预期。
  \item 当前公民对政府AI部署的开放态度可能不会持续，随着AI从抽象变为日常体验，谨慎案例和争议事件将塑造公众认知，政府需在窗口期内行动。
\end{itemize}
\begin{figure}[htbp]
\centering
\includegraphics[width=0.88\linewidth]{figures/fig_143.jpg}
\caption{EXHIBIT 2 - \#\# EXHIBIT 2 Citizens’ AI Proficiency Is Growing Fast, with Experts Almost Seven Times as Positive About AI’s Benefits as Those with No AI Experience Respondents' self-assessed level of AI knowledge (\%) Respondents who view the benefits of using AI in governme}
\end{figure}
\subsection{公民对代理型AI接受度超预期，偏好外部服务而非内部决策}
\textbf{公民对代理型AI（AI直接代表个人执行任务）的接受度超出预期，但偏好集中在外部服务场景（如聊天机器人）而非内部决策（如行政裁决），政府应优先部署面向公民的服务型AI，并建立“人在回路中”的监督机制。}
\begin{itemize}
  \item 公民最舒适的代理型AI场景是“7x24小时聊天机器人获取信息与帮助”，对内部决策类AI（如自动审批）接受度较低，表明信任集中在低风险、高透明度的服务。
  \item 报告建议政府优先部署面向公民的服务场景，如聊天机器人和虚拟助手，而非先做内部效率优化，以匹配公民偏好并快速建立信任。
  \item 建立“人在回路中”的监督机制，让公民知道AI决策可以被人类审查和推翻，是提升接受度的关键。
  \item 公民和公务员的AI素养培训至关重要，因为熟练度直接决定态度，政府应根据本地实际情况调整服务上线顺序和复杂度。
\end{itemize}
\begin{figure}[htbp]
\centering
\includegraphics[width=0.88\linewidth]{figures/fig_144.jpg}
\caption{Exhibit 3 - As this more autonomous version of AI emerged, industry observers assumed that users would be slow to adopt it, given agentic AI's increasingly immersive and imposing nature. But this has not been the case. In a short amount of time, citizens are demonstrating}
\end{figure}
\newpage
\section{辅助信号｜智库/国际机构}
本板块整合亚洲开发银行、国际清算银行、IMF和经合组织的最新研究，提炼出三个主题：能源系统气候适应缺口、汇率传导机制的结构性驱动因素、以及新兴市场与发展中经济体的融资与税制挑战。
\subsection{能源系统气候适应缺口：政策框架与资金错配}
\textbf{亚太地区能源基础设施正遭受极端天气冲击，但国家气候战略中能源适应优先级偏低，且适应资金流向与需求之间存在结构性错配，亟需系统性提升制度质量与融资工具多元化。}
\begin{itemize}
  \item 截至2024年，全球仅41\%的国家自主贡献将能源列为适应优先领域；国家适应计划中仅27\%明确纳入能源适应，亚洲开发银行发展中成员国表现略好（NDC 43\%、NAP 63\%），但整体仍严重不足。
  \item 2023年多边开发银行在能源、交通等建成环境领域的适应融资承诺仅72.5亿美元，占该部门气候融资总额（764.5亿美元）的不到10\%，私营部门参与意愿受前期成本高、长期收益低估、气候数据不确定性及政策不确定性制约。
  \item 气候信息服务可及性与准确性不足，制约了能源系统适应决策质量；报告识别出四大关键障碍：政策框架缺口、资金错配、数据短板和机构能力不足。
  \item 边际变化：极端天气频率上升正将能源适应从远期规划推为当前工程，但政策框架和资金流调整速度滞后于物理风险升级。
\end{itemize}
\begin{figure}[htbp]
\centering
\includegraphics[width=0.88\linewidth]{figures/fig_126.jpg}
\caption{Figure 1 - Figure 1: Exposed Value and Average Annual Loss by Infrastructure Sector Global Exposed Value (\textbackslash{}\$ billion) Average Annual Loss (\textbackslash{}\$ billion) Climate-driven disruptions in the energy sector are expected to intensify in the coming}
\end{figure}
\begin{figure}[htbp]
\centering
\includegraphics[width=0.88\linewidth]{figures/fig_128.jpg}
\caption{Figure 3 - Figure 4: Share of Adaptation Priorities in All NAPs as of 30 September 2023 NAP = national adaptation plan.}
\end{figure}
\subsection{汇率传导的结构性驱动：通胀与宏观规模超越贸易开放度}
\textbf{国际清算银行基于98个地区四十年数据，采用随机森林方法发现，通胀水平和宏观环境规模是解释汇率传导强度的最强变量，贸易开放度的重要性远低于理论预期，且发达地区传导系数长期下降。}
\begin{itemize}
  \item 发达地区一年期汇率传导系数从全样本的7.6\%降至2010年代的4.4\%，新兴市场与发展中地区维持在17\%-20\%之间，差距并非单纯由汇率制度或开放度决定。
  \item 随机森林分析显示，宏观环境规模与通胀水平是影响传导程度的最强因素：大型地区企业更倾向于定价到市场，降低传导；高通胀环境下企业更频繁调价，加速传导。
  \item 产品同质性与汇率波动性紧随其后成为第三、第四重要因素，同质性越高替代性越强，汇率变动越易通过竞争渠道传导至价格。
  \item 边际变化：发达地区传导系数持续下降，意味着在低通胀环境下，汇率波动对消费者价格的冲击可能进一步减弱，对货币政策传导和通胀预测有重要含义。
\end{itemize}
\begin{figure}[htbp]
\centering
\includegraphics[width=0.88\linewidth]{figures/fig_131.jpg}
\caption{Figure 1 - We classify all economies whose GDP in 2023 exceeded \textbackslash{}\$30,000 per capita as advanced economies. This group comprises 25 economies and includes countries whose income level is currently at or above those of Italy and South Korea. \$\textasciicircum{}\{8\}\$ The remaining countries }
\end{figure}
\begin{figure}[htbp]
\centering
\includegraphics[width=0.88\linewidth]{figures/fig_134.jpg}
\caption{Figure 4 - Fig. 3 - Mean Absolute Shapley Values by Factor (normalized) Importantly, we can go further in the analysis by eliciting the equivalent of partial effects in econometrics. As in the latter, these are the fitted pass-through models where all variables but one a}
\end{figure}
\subsection{新兴市场与发展中经济体的融资与税制挑战：从援助依赖到制度转型}
\textbf{IMF和经合组织分别指出，冈比亚税收缺口低于同类国家但制度质量成为天花板，东南亚蓝色经济面临2.1万亿美元融资缺口且过度依赖官方发展援助，两者均需系统性转向国内收入动员或多元化资本来源。}
\begin{itemize}
  \item 冈比亚税收收入占GDP约13\%，高于西非低收入国家均值，但低于15\%充足性门槛；约25\%来自关税和贸易相关税种，面临外部援助削减和非洲大陆自贸区推进的双重冲击。
  \item IMF随机前沿分析显示，冈比亚税收缺口约为GDP的2.5个百分点，低于低收入国家平均水平，意味着修补现有漏洞空间有限，下一步增长必须依赖制度质量提升和税基多元化。
  \item 东南亚蓝色经济到2030年面临近2.1万亿美元融资缺口（累计），当前海洋相关ODA年均仅约5.5亿美元，且DAC成员国对东南亚的海洋ODA在2024-2029年间可能持续下降。
  \item 边际变化：冈比亚从税种多元化转向税基和工具多元化，东南亚从ODA依赖转向系统性撬动私人资本和混合融资，两者均反映新兴市场融资模式的结构性转型压力。
\end{itemize}
\begin{figure}[htbp]
\centering
\includegraphics[width=0.88\linewidth]{figures/fig_136.jpg}
\caption{Figure 1 - \# 1 Blue economy trends and pressures in Southeast Asian countries Southeast Asia has emerged as the world's fourth-largest economic bloc, with its blue economy – economic activities tied to coastal, marine and freshwater resources – playing an important role }
\end{figure}
\begin{figure}[htbp]
\centering
\includegraphics[width=0.88\linewidth]{figures/fig_137.jpg}
\caption{Figure 1 - The data – available only for the wider Southeast Asia and Oceania region, rather than Southeast Asia specifically – show that the region has one of the highest relative contributions of the ocean economy globally. As shown in Figure 1.2, ocean economic activi}
\end{figure}
\section{结语}
后续需验证：韩国去杠杆后波动率能否正常化；中国结汇对人民币升值的持续推动力；CoWoS-L封装瓶颈的解决时间表；锂供给增量能否如期释放；欧洲天然气出口恢复速度；以及AI安全支出拐点是否提前至2026年底。
\newpage
\section{覆盖概览}
投行/券商 60 篇；战略咨询 1 篇；智库/国际机构 4 篇。\par
投行覆盖：GS 15篇 / JPM 11篇 / MS 11篇 / Citi 7篇 / JEF 6篇 / Bernstein 3篇 / HSBC 3篇 / DB 2篇 / BARC 1篇 / BofA 1篇\par
\section{Disclaimer}
{\scriptsize\color{refgray}
This community is a paid learning and information-sharing community created for general educational purposes only. The membership fee is charged solely for access to community discussions, educational content, organization, and operational costs. It is not an advisory fee, management fee, performance fee, brokerage fee, or compensation for personalized financial, investment, legal, tax, or professional advice.\par
I participate and share information solely as an individual and for educational discussion among community members. I am not acting as a financial adviser, investment adviser, broker, dealer, portfolio manager, legal adviser, tax adviser, accountant, fiduciary, or any other licensed professional adviser. Nothing shared in this community constitutes, and should not be interpreted as, financial advice, investment advice, legal advice, tax advice, a recommendation, solicitation, offer, endorsement, instruction, or guarantee to buy, sell, hold, short, trade, or invest in any security, cryptocurrency, fund, derivative, strategy, or other financial product.\par
All content shared in this community, including messages, discussions, links, files, charts, examples, market commentary, watchlists, AI-generated or AI-polished summaries, and third-party materials, is general, impersonal, and educational in nature. It does not take into account any individual's financial situation, investment objectives, risk tolerance, investment horizon, liquidity needs, tax status, legal restrictions, or suitability requirements. No content should be relied upon as suitable for any specific person, account, portfolio, or financial decision.\par
Information shared in this community may be incomplete, inaccurate, outdated, speculative, AI-generated, AI-edited, or based on third-party internet sources that have not been independently verified. I make no representation or warranty as to the accuracy, completeness, timeliness, reliability, or suitability of any information shared.\par
Financial markets involve substantial risk, including the possible loss of principal. Past performance, historical data, hypothetical examples, simulated results, backtests, personal opinions, or educational examples do not guarantee future results. Each member is solely responsible for conducting their own research, exercising independent judgment, and making their own decisions. Before making any financial, investment, legal, or tax decision, members should consult qualified and licensed professionals.\par
Participation in this community, payment of any membership fee, or communication with me or other members does not create any adviser-client, fiduciary, professional, contractual, agency, partnership, employment, or other advisory relationship. By joining or participating in this community, you acknowledge and agree that you are solely responsible for your own decisions, actions, transactions, profits, losses, risks, and consequences, and that I shall not be liable for any direct, indirect, incidental, consequential, financial, legal, tax, or other loss or damage arising from or related to any content, discussion, information, or material shared in this community.\par
}
\newpage
\begin{center}
{\Large 更多详情报告kcdesk.com}\par\vspace{0.8cm}
\includegraphics[width=0.58\linewidth]{\detokenize{../../prompts/zsxq_img.jpg}}
\end{center}
\end{document}
